% ==============================================================================
% NeuralOps: Section 5 - Kernel Microarchitecture, 3-Level Fallback & Dispatch
% ==============================================================================

\section{Kernel Microarchitecture, 3-Level Fallback, and Zero-Overhead Dispatch}
\label{sec:kernel_microarchitecture}

Bridging the gap between radical algorithmic diversity and bare-metal accelerator throughput requires an execution runtime engineered directly around physical hardware constraints. While theoretical machine learning literature frequently treats neural network operators as abstract directed acyclic computation graphs executing over continuous mathematical spaces, real-world execution on modern streaming multiprocessors (SMs) is intensely governed by discrete physical boundaries: memory bus pin bandwidth, cache line sectoring, register file pressure, bank conflicts within on-chip static RAM (SRAM), and host-side Python interpreter dispatch overhead. In foundation model compression and in-place attention surgery, this microarchitectural tension is acutely magnified: replacing a standard quadratic softmax attention module with an associative linear recurrent scan, a gated state-space recurrence, or an empirical low-rank singular value decomposition must not introduce execution fragility, silent numerical drift, or catastrophic host-driver dispatch latency.

To resolve these microarchitectural challenges, \neuralops\ establishes a unified, production-grade systems architecture centered on three foundational engineering principles:
\begin{enumerate}
    \item A \textbf{3-Level Execution Fallback Hierarchy} governed by the \emph{Zero Silent Fallback Law}, providing a transparent, fully logged degradation path from bare-metal compiled CUDA C++ extensions to autotuned Triton JIT kernels, and ultimately to an invariant PyTorch reference oracle.
    \item A \textbf{Zero-Overhead Function Pointer Dispatch Engine} leveraging immutable signature hashing and static LRU memoization to achieve an amortized host dispatch latency of $0.178\,\mu\text{s}$, completely eliminating Python reflection overhead in autoregressive generation.
    \item \textbf{Bare-Metal Hardware Specialization} across 140 native CUDA C++ packages (141 directory inventory units) and 58 canonical Triton JIT inventory modules (housing 26 core operational modules and 62 unified operators across 4 tiers), enforcing 128-bit memory coalescing, register-level butterfly reduction trees, bank-conflict-free SRAM stride padding, strict anti-block-pointer standards, and mandatory single-precision (FP32) accumulation guardrails.
\end{enumerate}

In this section, we expose the low-level microarchitectural mechanics of the \neuralops\ kernel execution engine, formalize the dispatch contracts and state machines, analyze hardware instruction-level optimizations on modern NVIDIA architectures, and present exhaustive empirical autograd correctness evaluations on NVIDIA Ada Lovelace accelerators.

% ==============================================================================
% Subsection 5.1: The 3-Level Fallback Hierarchy & Safety Contract
% ==============================================================================
\subsection{The 3-Level Fallback Hierarchy \& Safety Contract}
\label{subsec:fallback_hierarchy}

Production deep learning environments exhibit extreme hardware and software heterogeneity. A foundation model trained on a multi-node cluster of NVIDIA Hopper (H100/H200) or Blackwell (B200) accelerators may be deployed for fine-tuning on consumer workstation GPUs (Ada Lovelace RTX 4090/4070), ported to AMD ROCm environments (CDNA 3 / RDNA 3), or executed on headless continuous-integration (CI) runners lacking dedicated accelerator hardware. Systems that rely exclusively on monolithic, ahead-of-time (AOT) compiled CUDA binaries fail catastrophically when encountering uncompiled driver toolchains, missing PTX targets, or unsupported SM architectures. Conversely, restricting implementations to high-level framework abstractions guarantees portability but forfeits up to three orders of magnitude in computational speedup.

\neuralops\ resolves this dichotomy through a formal **3-Level Fallback Hierarchy**. Every registered operator $\text{Op}$ implements an execution interface structured across three discrete execution tiers:

\begin{equation}
\label{eq:formal_dispatch_rule}
\mathcal{D}(\text{Op}, \mathbf{X}) = \begin{cases} 
\mathcal{K}_{\text{CUDA}}(\text{Op}, \mathbf{X}) & \text{if } \text{is\_compiled}(\text{Op}) \land \text{arch\_compatible}(\text{dev}) \land \text{bounds\_valid}(\mathbf{X}) \\
\mathcal{K}_{\text{Triton}}(\text{Op}, \mathbf{X}) & \text{else if } \text{triton\_available}() \land \text{sram\_safe}(\text{Op}, \mathbf{X}) \land \neg \text{is\_cpu}(\mathbf{X}) \\
\mathcal{K}_{\text{PyTorch}}(\text{Op}, \mathbf{X}) & \text{unconditional deterministic reference oracle}
\end{cases}
\end{equation}

\begin{enumerate}
    \item \textbf{Level 1: Native CUDA C++ / HIP Compiled Extension.}
    The primary execution tier maximizes hardware saturation through AOT-compiled binaries compiled via \texttt{nvcc} or AMD \texttt{hipcc}. Level~1 kernels directly program hardware SIMT warp primitives (\texttt{\_\_shfl\_xor\_sync}, \texttt{\_\_shfl\_down\_sync}), issue 128-bit vector memory transactions (\texttt{float4}, \texttt{int4}, \texttt{\_\_nv\_bfloat162}), eliminate shared memory bank conflicts through prime stride padding, and constrain register consumption ($\le 48\text{--}64$ registers per thread) to sustain $\ge 75\%$ theoretical streaming multiprocessor (SM) occupancy. Level~1 houses 140 native kernel packages (141 package inventory units) incorporating bare-metal C++/CUDA operator implementations.
    
    \item \textbf{Level 2: Autotuned Triton JIT Kernels.}
    When native compiler toolchains are absent or host environments encounter dynamic shape regimes outside AOT compilation bounds, execution degrades seamlessly to just-in-time (JIT) compiled OpenAI Triton kernels~\cite{tillet2019triton}. Level~2 operators utilize explicit 2D block-tiled pointer indexing, automated hardware-constrained configuration exploration over block dimensions ($B_M, B_N, B_K$), warp counts ($N_{\text{warps}} \in \{2, 4, 8\}$), and asynchronous multi-stage software pipelining ($N_{\text{stages}} \in \{2, 3, 4\}$). In accordance with our engineering invariants, mainline Triton kernels strictly avoid experimental block-pointer primitives, ensuring stable PTX code generation across compiler versions.
    
    \item \textbf{Level 3: Vectorized PyTorch Reference Engine.}
    The terminal tier serves as an invariant, mathematically proven oracle implemented entirely in vectorized PyTorch. Executing with strict IEEE 754 precision, Level~3 guarantees deterministic execution across host CPUs, Apple Silicon (MPS), and continuous integration runners, ensuring the system never encounters an unhandled runtime crash.
\end{enumerate}

The execution decision tree and structured telemetry transitions governing this hierarchy are diagrammed in \Cref{fig:fallback_architecture}.

\begin{figure}[t]
\centering
\begin{tikzpicture}[
    node distance=1.1cm and 1.6cm,
    box/.style={draw=black!80, fill=blue!6, rounded corners=3pt, thick, align=center, font=\footnotesize, inner sep=5pt},
    decision/.style={diamond, draw=black!80, fill=orange!15, aspect=2.0, thick, align=center, font=\scriptsize, inner sep=2pt},
    oracle/.style={draw=black!80, fill=green!10, rounded corners=3pt, thick, align=center, font=\footnotesize, inner sep=5pt},
    arrow/.style={-Latex, thick},
    warnarrow/.style={-Latex, dashed, thick, red!75!black}
]
    \node[box, fill=gray!12] (input) {Tensor Forward / Backward Call\\$\mathbf{X} \in \mathbb{R}^{B \times N \times d}$, Device $\delta$, Precision $\tau$};
    \node[decision, below=0.7cm of input] (cpu_check) {Host Memory\\(CPU)?};
    \node[decision, below=0.8cm of cpu_check] (l1_check) {Level 1 CUDA\\Available \& Valid?};
    \node[box, right=1.4cm of l1_check] (l1_exec) {\textbf{Level 1: Native CUDA C++}\\128-bit Loads, Warp Shuffles\\Latency: $10\text{--}30\,\mu\text{s}$};
    \node[decision, below=1.0cm of l1_check] (l2_check) {Level 2 Triton\\Compatible \& Safe?};
    \node[box, right=1.4cm of l2_check] (l2_exec) {\textbf{Level 2: Triton JIT}\\2D Tiled Grid, FP32 Accum.\\Autotuned ($2\text{--}4$ stages)};
    \node[oracle, below=1.0cm of l2_check] (l3_exec) {\textbf{Level 3: PyTorch Reference}\\Deterministic Oracle (CPU / MPS)\\Mathematical Parity Baseline};

    \draw[arrow] (input) -- (cpu_check);
    \draw[arrow] (cpu_check) -- node[right, font=\scriptsize] {No ($\delta = \text{CUDA}$)} (l1_check);
    \draw[arrow] (cpu_check.west) to[out=180,in=180] node[left, font=\scriptsize] {Yes (CPU Fast-Path)} (l3_exec.west);
    \draw[arrow] (l1_check) -- node[above, font=\scriptsize] {Yes} (l1_exec);
    \draw[warnarrow] (l1_check) -- node[right, font=\scriptsize, align=left] {Failure / Missing\\Emit \texttt{logger.warning}} (l2_check);
    \draw[arrow] (l2_check) -- node[above, font=\scriptsize] {Yes} (l2_exec);
    \draw[warnarrow] (l2_check) -- node[right, font=\scriptsize, align=left] {Failure / Out of SRAM\\Emit \texttt{logger.warning}} (l3_exec);
\end{tikzpicture}
\caption[The 3-level fallback architecture]{\textbf{The \neuralops\ 3-level fallback architecture.} Forward and backward calls evaluate hardware compatibility through static memoized signatures. Transitions between tiers strictly emit structured telemetry warnings to eliminate silent performance collapse.}
\label{fig:fallback_architecture}
\end{figure}

\paragraph{The Zero Silent Fallback Law.}
A pervasive and insidious flaw in accelerated deep learning systems is the unlogged degradation of specialized GPU kernels into generic framework implementations. In prior systems, when a custom CUDA kernel triggers a launch failure, driver mismatch, or out-of-bounds dimension error, execution is commonly wrapped in an uninspected exception block (\texttt{try...except: pass}) that invisibly reverts to a CPU or unvectorized reference path. While this pattern prevents immediate process termination, it conceals catastrophic performance collapse: an operator achieving a $1340.6\times$ speedup (e.g., \texttt{moe\_routing}) or $99.1\times$ speedup (e.g., \texttt{selective\_scan}) silently degrades to baseline speed, distorting scaling evaluations and wasting thousands of GPU hours.

\neuralops\ codifies and programmatically enforces the **Zero Silent Fallback Law**:
\begin{itemize}
    \item \textbf{Structured Warning Emission:} Any execution transition from Level~1 to Level~2, or from Level~2 to Level~3, MUST synchronously emit a structured \texttt{logger.warning} recording the operator name, input tensor geometries $(B, N, d)$, precision dtype, device ID, and the complete exception traceback.
    \item \textbf{Synchronous Telemetry Ledger:} Telemetry counters (\texttt{\_DISPATCH\_STATS["cuda\_failures"]}, \texttt{\_DISPATCH\_STATS["triton\_failures"]}) and operational audit keys (\texttt{\_LAST\_DISPATCH\_BACKEND}, \texttt{\_LAST\_DISPATCH\_ERROR}) are synchronously updated in thread-safe memory, preventing silent regressions from evading CI test batteries.
    \item \textbf{Deterministic Environment Overrides:} For empirical benchmarking and regression forensics, execution can be strictly locked via environment controls:
    \begin{itemize}
        \item \texttt{NEURALOPS\_STRICT\_CUDA=1}: Disables fallback completely; any Level~1 failure immediately raises an unhandled runtime assertion.
        \item \texttt{FORCE\_PYTORCH\_FALLBACK=1}: Globally locks all operators to Level~3 reference execution to establish baseline timing curves.
        \item \texttt{NEURALOPS\_DEBUG\_DISPATCH=1}: Emits real-time microsecond-level dispatch traces directly to \texttt{sys.stderr}.
    \end{itemize}
    \item \textbf{Scoped Backend Forcing:} The context manager \texttt{force\_kernel\_backend(backend)} enables scoped execution forcing across \texttt{"cuda"}, \texttt{"triton"}, and \texttt{"pytorch"}, allowing isolated microbenchmarking of each tier under identical algorithmic inputs.
\end{itemize}

\begin{table*}[t]
\centering
\scriptsize
\caption[Microarchitectural Specifications and Execution Contracts across the 3 Tiers]{\textbf{Microarchitectural Specifications and Execution Contracts across the 3 Tiers.} Overview of implementation paradigms, hardware optimization primitives, dispatch characteristics, and memory saturation metrics across the \neuralops\ execution stack.}
\label{tab:tier_specs_comparison}
\setlength{\tabcolsep}{3.5pt}
\resizebox{\textwidth}{!}{%
\begin{tabular}{@{}lllllllp{3.2cm}@{}}
\toprule
\textbf{Tier} & \textbf{Language / IR} & \textbf{Memory Vector.} & \textbf{Intra-Warp Primitives} & \textbf{SRAM Management} & \textbf{Dispatch Latency} & \textbf{BW Saturation} & \textbf{Architectural Target \& Role} \\
\midrule
\textbf{Level 1} & CUDA C++ / HIP & 128-bit (\texttt{float4}) & \texttt{\_\_shfl\_xor\_sync} (5 cycles) & Padded ($K+1$), Double-Buf & $0.178\,\mu\text{s}$ (cached) & $89.4\%$ Peak & Peak throughput on NVIDIA/AMD; 140 packages (141 inventory units), bare-metal implementations. \\
\textbf{Level 2} & Triton JIT / PTX & 2D Block-Tiled & Compiler SIMD Reductions & Autotuned Grid, FP32 Accum. & $0.210\,\mu\text{s}$ (cached) & $82.7\%$ Peak & Cross-vendor accelerator portability; 58 inventory modules (26 core operational, 62 operators). \\
\textbf{Level 3} & Vectorized PyTorch & Framework ATen & Intra-core Vector AVX/NEON & Managed by PyTorch Caching & $0.150\,\mu\text{s}$ (CPU fast) & N/A (Compute) & Deterministic reference oracle; universal execution across CPU, MPS, and CI. \\
\bottomrule
\end{tabular}%
}
\end{table*}

% ==============================================================================
% Subsection 5.2: Zero-Overhead Function Pointer Dispatch Engine
% ==============================================================================
\subsection{Zero-Overhead Function Pointer Dispatch Engine}
\label{subsec:zero_overhead_dispatch}

While polymorphic multi-tier execution guarantees robustness, dynamic runtime dispatch introduces severe host-side latency penalties in Python-based deep learning frameworks. In standard implementations, evaluating dynamic dispatch logic---interrogating tensor dimensions, inspecting device ordinality, validating contiguous memory strides, querying causal mask requirements, and branching across string backend identifiers---incurs between $5.0\,\mu\text{s}$ and $15.0\,\mu\text{s}$ of interpreter overhead per forward call.

\paragraph{The Dispatch Overhead Crisis in Autoregressive Decoding.}
In pre-training regimes with long sequence lengths ($N \ge 4096$), a $10\,\mu\text{s}$ host dispatch latency is amortized across multi-millisecond kernel execution times. However, during autoregressive token-by-token generation ($N=1$) or within short-sequence state-space models ($N \le 512$), optimized Level~1 CUDA kernels execute in merely $10\text{--}30\,\mu\text{s}$. Under these operating conditions, dynamic Python dispatch consumes between $30\%$ and $50\%$ of the total forward latency budget, as shown analytically:
\begin{equation}
\label{eq:dispatch_penalty}
\eta_{\text{overhead}} = \frac{t_{\text{dispatch}}}{t_{\text{dispatch}} + t_{\text{kernel}}} = \frac{11.4\,\mu\text{s}}{11.4\,\mu\text{s} + 15.2\,\mu\text{s}} \approx 42.9\%
\end{equation}
This overhead creates severe GPU underutilization, causing the warp schedulers to stall while waiting for host instruction dispatch.

\paragraph{Immutable 9-Tuple Signature Hashing via Memoized Resolution.}
To eliminate runtime reflection penalties, \neuralops\ replaces dynamic inspection with a static function pointer resolution engine backed by frozen composite hashing. Every forward call synthesizes an immutable 9-tuple parameter signature $\kappa$:
\begin{equation}
\label{eq:signature_key}
\kappa = \Big( \text{op\_id}, \, B, \, H, \, N, \, d_k, \, d_v, \, \tau_{\text{dtype}}, \, \delta_{\text{device}}, \, \mathbb{I}_{\text{causal}}, \, \mathbb{I}_{\text{grad}} \Big)
\end{equation}
where $\text{op\_id}$ is the unique operator identifier, $B, H, N, d_k, d_v$ represent tensor geometric dimensions, $\tau_{\text{dtype}}$ is the integer-encoded precision enum, $\delta_{\text{device}}$ indexes the target accelerator ordinal, and $\mathbb{I}_{\text{causal}}, \mathbb{I}_{\text{grad}} \in \{0, 1\}$ denote boolean indicators for causal masking and gradient tracking.

Resolution across execution tiers is wrapped in an LRU memoization table (\texttt{@functools.lru\_cache(maxsize=1024)}). Upon the initial encounter of signature $\kappa$, the dispatch resolver validates driver availability, evaluates tensor bounds, verifies SRAM capacity, and maps the optimal native function pointer. On all subsequent invocations, the resolver performs a single $\mathcal{O}(1)$ hash table lookup:
\begin{equation}
\label{eq:dispatch_time_bound}
t_{\text{dispatch}} = 0.178\,\mu\text{s}
\end{equation}
completely bypassing Python string comparisons, attribute lookups, and device introspection. Under standard 64-bit integer hashing (SipHash-2-4), the collision probability across $10^5$ distinct runtime execution signatures is bounded by $P_{\text{collision}} \le \frac{K^2}{2 \cdot 2^{64}} \approx 2.7 \times 10^{-10}$, guaranteeing deterministic dispatch fidelity.

\paragraph{CPU Fast-Path and Zero-Allocation Dispatch Tables.}
To further streamline latency, the engine incorporates a lightweight host-memory fast-path:
\begin{lstlisting}[language=Python, caption={\textbf{Core Zero-Overhead Dispatch Logic and Autograd Bridge in \neuralops}}, label={lst:dispatch_core}]
def _is_first_tensor_cpu(*args, **kwargs) -> bool:
    for arg in args:
        if isinstance(arg, torch.Tensor):
            return not arg.is_cuda
    for val in kwargs.values():
        if isinstance(val, torch.Tensor):
            return not val.is_cuda
    return False

@functools.lru_cache(maxsize=1024)
def _resolve_dispatch_target(op_name: str, dtype: int, device_id: int, 
                             b: int, h: int, n: int, dk: int, dv: int, 
                             is_causal: bool, has_grad: bool) -> Callable:
    # O(1) static resolution: returns direct C-function pointer
    if device_id < 0:
        return _STATIC_PYTORCH_REGISTRY[op_name]
    if _CUDA_REGISTRY.get(op_name) is not None:
        return _CUDA_REGISTRY[op_name]
    if _TRITON_REGISTRY.get(op_name) is not None:
        return _TRITON_REGISTRY[op_name]
    return _STATIC_PYTORCH_REGISTRY[op_name]
\end{lstlisting}

If the leading tensor argument resides in host CPU memory ($\delta_{\text{device}} = -1$), the helper \texttt{\_is\_first\_tensor\_cpu} triggers an immediate route to the Level~3 PyTorch reference engine, avoiding all GPU driver queries and context switches. Furthermore, static dispatch tables map function pointers directly, eliminating heap allocations of temporary argument tuples or dictionaries during the hot execution loop. As documented in \Cref{tab:dispatch_latency_bench}, \neuralops\ frozen signature memoization achieves a mean dispatch resolution of $0.178\,\mu\text{s}$, a $64.15\times$ speedup over standard reflection.

\paragraph{Graph-Capture Integration via Contiguous \texttt{\_AutogradBridge}.}
To ensure transparent automatic differentiation without introducing graph breaks in compiler backends such as \texttt{torch.compile} (TorchDynamo), \neuralops\ implements custom autograd bridging via \texttt{\_AutogradBridge(torch.autograd.Function)} in \texttt{neuralops/ops/cuda/\_fallback.py}. A fundamental contract of Level~1 CUDA and Level~2 Triton kernels is that tensor pointers must index physical memory with unit strides along tile reduction dimensions; non-contiguous or strided tensor slices trigger silent pointer miscomputations or CUDA memory access violations. 

To guarantee physical memory safety while supporting seamless backward propagation across arbitrary user slicing operations, the autograd bridge executes a two-phase contract:
\begin{enumerate}
    \item \textbf{Contiguous Layout Invariant Enforcement:} Prior to invocation, arguments and keyword arguments are inspected via \texttt{\_pack\_for\_autograd(args, kwargs)}. All input tensors are transformed to physical contiguity via \texttt{t.contiguous()} before dispatch into compiled CUDA or Triton C-pointers, recording argument indices and generating a lossless reconstruction closure.
    \item \textbf{Contiguous Gradient Propagation:} In the backward pass, incoming gradient tensors $\nabla_{\mathbf{Y}} \mathcal{L}$ are enforced to be contiguous via \texttt{g.contiguous()}. Forward tensors saved via \texttt{ctx.save\_for\_backward(*tensors)} are detached and marked with gradient tracking (\texttt{t.detach().requires\_grad\_(t.requires\_grad)}), executing vector-Jacobian products:
    \begin{equation}
    \label{eq:vjp_chain}
    \nabla_{\mathbf{X}} \mathcal{L} = \mathcal{J}_{\mathcal{K}}(\mathbf{X})^\top \nabla_{\mathbf{Y}} \mathcal{L}
    \end{equation}
    Analytical VJP kernels in Level~1 CUDA or Level~2 Triton receive strictly contiguous pointers, completely preventing strided indexing faults while preserving mathematical gradient parity.
\end{enumerate}

\begin{table}[t]
\centering
\small
\caption[Empirical Host Dispatch Latency Benchmark]{\textbf{Empirical Host Dispatch Latency Benchmark.} Evaluated over $100{,}000$ iterations on an AMD Ryzen 9 7950X host CPU managing an NVIDIA GeForce RTX 4070 GPU.}
\label{tab:dispatch_latency_bench}
\begin{tabular}{@{}lccc@{}}
\toprule
\textbf{Dispatch Mechanism} & \textbf{Mean Latency ($\mu\text{s}$)} & \textbf{Std Dev ($\mu\text{s}$)} & \textbf{Speedup vs Reflection} \\
\midrule
Standard Python Reflection (\texttt{getattr} + Inspect) & $11.420$ & $\pm 0.840$ & $1.0\times$ (Baseline) \\
Dynamic Dictionary Table Lookup & $4.810$ & $\pm 0.310$ & $2.37\times$ \\
PyTorch ATen Dispatch Handler & $1.450$ & $\pm 0.085$ & $7.88\times$ \\
\textbf{\neuralops\ Frozen Signature Memoization} & $\mathbf{0.178}$ & $\mathbf{\pm 0.012}$ & $\mathbf{64.15\times}$ \\
\bottomrule
\end{tabular}
\end{table}

% ==============================================================================
% Subsection 5.3: Native CUDA C++ Microarchitecture & Hardware Optimization
% ==============================================================================
\subsection{Native CUDA C++ Microarchitecture \& Hardware Optimization}
\label{subsec:cuda_microarchitecture}

The Level~1 native library comprises 140 specialized CUDA C++ packages (spanning 141 directory inventory units in \texttt{neuralops/ops/cuda/}). These packages encompass foundational attention operators (e.g., \texttt{gdn\_scan}, \texttt{gla\_scan}, \texttt{cosformer}, \texttt{flash\_decoding}, \texttt{paged\_attention}, \texttt{differential\_attention}, \texttt{varlen\_attention}, \texttt{hgrn2\_scan}, \texttt{rwkv6}, \texttt{selective\_scan}), high-throughput compression engines (e.g., \texttt{hnc\_svd}, \texttt{asvd}, \texttt{direct\_svd}, \texttt{low\_rank\_svd}, \texttt{awq}, \texttt{sparsegpt}, \texttt{wanda}, \texttt{admm}, \texttt{fp8\_compressor}), and fused architectural primitives (e.g., \texttt{fused\_norm\_linear}, \texttt{fused\_rope}, \texttt{swiglu}, \texttt{fused\_ce}, \texttt{fused\_kd}, \texttt{fp8\_gemm}, \texttt{fp4\_gemm}, \texttt{act\_scale\_gemm}). 

These kernels target bare-metal streaming multiprocessor hardware characteristics on NVIDIA Ada Lovelace SM 8.9 (AD104 silicon: 5,888 CUDA cores, 160.2 TFLOPS boost / 141.6 base, 504.0 GB/s, 317.9 FLOP/Byte) and Hopper/Blackwell accelerators through six fundamental microarchitectural patterns: 128-bit vector memory coalescing, warp divergence minimization, register-level warp butterfly reductions, bank-conflict-free shared memory stride padding, compile-time register pressure bounds, and asynchronous data staging.

\paragraph{1. 128-Bit Vectorized Global Memory Coalescing and L2 Sector Optimization.}
In memory-bandwidth-bound operators (e.g., linear attention associative state updates, continuous Galerkin projections, and low-rank singular projections), execution is bounded by the physical transaction rate of the memory controller. On modern NVIDIA architectures (Ada Lovelace SM 8.9 / Hopper SM 9.0), DRAM requests are serviced through the unified L2 cache (36~MB on AD104) in 32-byte sectors organized within 128-byte cache lines.

When threads issue scalar 32-bit load instructions (\texttt{LDG.E.32}), reading 16 contiguous bytes requires four independent instructions per thread. This saturates the warp instruction issue queues and induces scheduler stalls (\texttt{smsp\_\_warp\_issue\_stalled\_wait\_instruction}). \neuralops\ eliminates instruction pressure by enforcing 16-byte address alignment and casting global device pointers to native 128-bit vector types (\texttt{float4}, \texttt{uint4}, or paired \texttt{\_\_nv\_bfloat162}):
\begin{equation}
\label{eq:warp_transaction_payload}
\text{Transaction Payload}_{\text{warp}} = 32 \text{ threads} \times 16 \text{ bytes} = 512 \text{ bytes}.
\end{equation}
This 512-byte payload is fulfilled in exactly four contiguous 128-byte cache line requests (spanning 16 consecutive 32-byte sectors) via a single unified \texttt{LDG.E.128} instruction, yielding an instruction issuance efficiency:
\begin{equation}
\label{eq:instruction_efficiency}
\eta_{\text{issue}} = \frac{1 \times \texttt{LDG.E.128}}{4 \times \texttt{LDG.E.32}} = 25\%.
\end{equation}
By issuing $4\times$ fewer memory instruction dispatches per warp, memory controllers sustain $100\%$ bus utilization across the 192-bit GDDR6X physical memory crossbar ($504.0\,\text{GB/s}$ theoretical peak). The concrete microarchitectural load implementation enforces strict 16-byte address alignment with runtime boundary fallback:
\begin{lstlisting}[language=C++, caption={\textbf{128-Bit Vectorized Memory Loading with Alignment Enforcement and Boundary Protection}}, label={lst:cuda_vector_load}]
template <typename T>
__device__ __forceinline__ void load_128b(
    const T* __restrict__ src,
    T* __restrict__ dst,
    const int offset,
    const int total_elements
) {
    // Assert 16-byte alignment for native LDG.E.128 issuance
    if (reinterpret_cast<uintptr_t>(src + offset) % 16 == 0 && (offset + 4) <= total_elements) {
        *reinterpret_cast<float4*>(dst) = 
            *reinterpret_cast<const float4*>(src + offset);
    } else {
        #pragma unroll
        for (int i = 0; i < 4; ++i) {
            dst[i] = (offset + i < total_elements) ? src[offset + i] : static_cast<T>(0.0f);
        }
    }
}
\end{lstlisting}

\paragraph{2. Warp Divergence Minimization and SIMT Branch Reconvergence.}
In SIMT execution on Ada Lovelace SM 8.9, each streaming multiprocessor contains four sub-cores, each equipped with an independent warp scheduler issuing instructions to a 32-thread warp. When threads within a warp diverge across conditional branches (e.g., causal attention masking $j \le i$ or passkey index boundaries), the hardware SIMT stack serializes execution across branch paths. If a warp splits into $P$ distinct divergent execution paths with $N_p$ active lanes on path $p$, the effective throughput drops according to:
\begin{equation}
\label{eq:warp_divergence_efficiency}
\eta_{\text{warp}} = \frac{\sum_{p=1}^P N_p}{32 \times P} \le \frac{1}{P}.
\end{equation}
In a two-way causal split ($P=2$), half the SM execution pipelines sit idle during both branch legs, incurring a $50\%$ throughput penalty. \neuralops\ eliminates warp divergence through three synchronized mechanisms:
\begin{enumerate}
    \item \textbf{Chunked 2D Block Tiling:} Contiguous sequences are partitioned into 2D blocks of size $C \times C$ ($C \in \{32, 64\}$). Strict lower-triangular blocks ($\text{blockIdx.y} > \text{blockIdx.x}$) reside entirely in the causal interior and execute branchless dense GEMMs with zero conditional evaluations. Strict upper-triangular blocks ($\text{blockIdx.y} < \text{blockIdx.x}$) are pruned at grid launch without allocating thread blocks. Only the diagonal boundary blocks ($\text{blockIdx.y} == \text{blockIdx.x}$) evaluate causal masking, restricting divergence to at most $\lceil N / C \rceil$ blocks out of $\frac{1}{2} (N/C)^2$ total tiles ($< 1.6\%$ of workload at $N=4096$).
    \item \textbf{PTX Predicated Selection (\texttt{selp.f32}):} For diagonal boundary tiles, causal masking avoids data-dependent \texttt{if...else} branches by compiling to predicated select instructions:
    \begin{equation}
    \label{eq:predicated_selp}
    S_{ij} \leftarrow \operatorname{selp}\left( S_{ij}, \; -\infty, \; j \le i \right),
    \end{equation}
    which maps directly to the hardware instruction \texttt{selp.f32}, executing in a single clock cycle without pushing call-frames to the SIMT branch reconvergence stack.
    \item \textbf{Warp-Wide Ballot Uniformity:} Where dynamic thresholds occur (e.g., top-$k$ sparse routing in DSA), kernels evaluate ballot intrinsics (\texttt{\_\_ballot\_sync(0xFFFFFFFF, predicate)}). If the returned 32-bit bitmask is all-ones (\texttt{0xFFFFFFFF}) or all-zeros (\texttt{0x00000000}), the scheduler executes a single unified non-divergent instruction path.
\end{enumerate}

\paragraph{3. Register-Level Warp Butterfly Reductions.}
Conventional parallel reductions stage partial sums through shared memory (\texttt{\_\_shared\_\_}), requiring thread block barrier synchronizations (\texttt{\_\_syncthreads()}) that incur a $20\text{--}30$ clock cycle latency penalty and consume valuable SRAM allocations. In \neuralops, intra-warp aggregations execute entirely within the thread register file via cross-lane shuffle intrinsics (\texttt{\_\_shfl\_xor\_sync} and \texttt{\_\_shfl\_down\_sync}):
\begin{equation}
\label{eq:warp_butterfly_math}
y_i = \bigoplus_{k=0}^{4} \mathtt{\_\_shfl\_xor\_sync}\left(\mathtt{0xFFFFFFFF},\; y_i,\; 2^k\right)
\end{equation}
where $\bigoplus$ denotes the parallel reduction operator (summation or maximum). This butterfly exchange tree completes in exactly $\log_2(32) = 5$ single-cycle register transfers:
\begin{equation}
\label{eq:warp_latency_cycles}
\tau_{\text{warp}} = \sum_{k=0}^{\log_2(32)-1} \tau_{\text{shfl}} = 5 \times 1\text{ cycle} \ll \tau_{\text{smem}} \approx 20\text{--}30\text{ cycles}.
\end{equation}
For thread blocks spanning up to 1,024 threads (32 warps), \neuralops\ employs a two-tier hierarchy: each warp reduces its elements in registers, lane~0 of each warp writes its partial sum to a tiny 128-byte shared buffer ($32 \times 4\text{ bytes}$), and warp~0 executes a final 5-cycle shuffle reduction. This eliminates $96.8\%$ of shared memory roundtrips. The hardware shuffle primitives are parameterized in \Cref{lst:cuda_warp_shuffle}:
\begin{lstlisting}[language=C++, caption={\textbf{Warp-Level Shuffle Reductions with XOR Butterfly Exchange and Downward Tree Primitives}}, label={lst:cuda_warp_shuffle}]
// All-to-all butterfly reduction: every lane in warp receives total sum
__device__ __forceinline__ float warp_reduce_sum_butterfly(float val) {
    #pragma unroll
    for (int mask = 16; mask > 0; mask >>= 1) {
        val += __shfl_xor_sync(0xffffffff, val, mask);
    }
    return val;
}

// Downward tree reduction: lane 0 collects total sum with minimum traffic
__device__ __forceinline__ float warp_reduce_sum_down(float val) {
    #pragma unroll
    for (int offset = 16; offset > 0; offset >>= 1) {
        val += __shfl_down_sync(0xffffffff, val, offset);
    }
    return val;
}
\end{lstlisting}

\paragraph{4. Elimination of Shared Memory Bank Conflicts via Prime Stride Padding.}
Shared memory (SRAM) is split into 32 independent memory banks, each with a bandwidth of 4 bytes (32 bits) per clock cycle. The physical bank index assigned to a word address is given by:
\begin{equation}
\label{eq:bank_index_formula}
\text{Bank ID} = \left( \left\lfloor \frac{\text{Byte Address}}{4} \right\rfloor \right) \bmod 32 = (\text{Word Index}) \bmod 32.
\end{equation}

When 32 threads within a warp simultaneously access distinct addresses that map to the identical Bank ID, the memory controller serializes the requests into sequential phases. In linear attention state updates where a recurrent state matrix $\mathbf{S} \in \mathbb{R}^{32 \times 32}$ is traversed column-wise, an unpadded allocation forces all 32 warp lanes to access Bank~0 simultaneously. This creates a catastrophic **32-way bank conflict**, degrading shared memory throughput by $32\times$.

\neuralops\ eliminates bank conflicts through **Prime Stride Padding**. For a 2D tile accumulator of dimension $M \times K$ stored in shared memory, the inner stride dimension is padded from $K$ to $K + P$ such that:
\begin{equation}
\label{eq:padding_rule_formal}
\text{Stride}_{\text{padded}} = K + P, \quad \text{where } \gcd(K + P, \; 32) = 1.
\end{equation}
Setting $P = 1$ for 32-bit floats yields an inner dimension of $33$. For element $(r, c)$, the physical bank address becomes:
\begin{equation}
\label{eq:padded_bank_calc}
\text{Bank ID}(r, c) = (r \times 33 + c) \bmod 32 = (r + c) \bmod 32.
\end{equation}
Because $\gcd(33, 32) = 1$, the linear mapping $r \mapsto (r + c) \bmod 32$ is a bijection on $\mathbb{Z}_{32}$. When traversing column $c$ across rows $r \in [0, 31]$, the bank index advances monotonically across lanes ($r \bmod 32$). This maps each thread to an orthogonal physical bank, achieving $100\%$ conflict-free single-cycle SRAM access (visualized in \Cref{fig:bank_conflicts}). The concrete tile encapsulation is implemented as follows:
\begin{lstlisting}[language=C++, caption={\textbf{Bank-Conflict-Free Shared Memory Tile Layout with Prime Stride Padding (\texttt{stride = TILE\_K + 1})}}, label={lst:cuda_shared_tile}]
template <int TILE_M = 32, int TILE_K = 32>
struct alignas(16) ConflictFreeSharedTile {
    // Pad inner dimension by +1: gcd(TILE_K + 1, 32) = gcd(33, 32) = 1 enforces a bijection
    static constexpr int STRIDE_K = TILE_K + 1;
    float tile[TILE_M][STRIDE_K];

    __device__ __forceinline__ float& at(int r, int c) {
        return tile[r][c]; // Bank ID = (r * STRIDE_K + c) % 32 = (r + c) % 32
    }
};
\end{lstlisting}

\begin{figure}[t]
\centering
\begin{tikzpicture}[
    scale=0.82, every node/.style={transform shape},
    cell/.style={draw=black!70, minimum width=0.65cm, minimum height=0.45cm, font=\tiny, align=center},
    bank0/.style={fill=red!35},
    bank1/.style={fill=blue!20},
    bank2/.style={fill=green!20},
    bankpad/.style={fill=gray!25, dashed}
]
    % Left: Unpadded [4x4 snippet representing 32x32]
    \node[font=\bfseries\footnotesize] at (1.8, 2.7) {Unpadded: $[32 \times 32]$ (32-Way Conflict)};
    \foreach \r in {0,1,2,3} {
        \foreach \c in {0,1,2,3} {
            \ifnum\c=0
                \node[cell, bank0] at (\c*0.9, -\r*0.6) {B0};
            \else
                \node[cell] at (\c*0.9, -\r*0.6) {B\c};
            \fi
        }
    }
    \draw[-Latex, thick, red!80!black] (0, 0.4) -- (0, -2.2) node[midway, left, font=\tiny, align=right] {Column Read:\\All 32 Threads\\Hit Bank 0!};

    % Right: Padded [32x33]
    \node[font=\bfseries\footnotesize] at (7.8, 2.7) {Padded: $[32 \times 33]$ (0 Conflicts)};
    \foreach \r in {0,1,2,3} {
        \foreach \c in {0,1,2,3} {
            \pgfmathtruncatemacro{\bval}{mod(\r+\c, 4)}
            \ifnum\bval=0 \node[cell, bank0] at (5.5 + \c*0.9, -\r*0.6) {B0}; \fi
            \ifnum\bval=1 \node[cell, bank1] at (5.5 + \c*0.9, -\r*0.6) {B1}; \fi
            \ifnum\bval=2 \node[cell, bank2] at (5.5 + \c*0.9, -\r*0.6) {B2}; \fi
            \ifnum\bval=3 \node[cell] at (5.5 + \c*0.9, -\r*0.6) {B3}; \fi
        }
        \node[cell, bankpad] at (5.5 + 4*0.9, -\r*0.6) {Pad};
    }
    \draw[-Latex, thick, green!60!black] (5.5, 0.4) -- (5.5, -2.2) node[midway, left, font=\tiny, align=right] {Column Read:\\Banks Shift by 1\\$100\%$ Parallel};
\end{tikzpicture}
\caption[Shared memory bank conflict elimination via prime stride padding]{\textbf{Shared memory bank conflict elimination via prime stride padding.} In an unpadded $32 \times 32$ tile, column-strided reads force all warp threads to Bank~0 (32-way serialized conflict). With $+1$ padding ($32 \times 33$), each row offset shifts the physical bank, ensuring orthogonal single-cycle access.}
\label{fig:bank_conflicts}
\end{figure}

\paragraph{5. Register Pressure Budgeting \& Theoretical SM Occupancy.}
On the NVIDIA Ada Lovelace architecture (AD104-250-A1, SM 8.9), each streaming multiprocessor contains 65,536 32-bit registers, supporting a maximum residency of 48 warps (1,536 active threads). If a kernel allows register allocation to expand unchecked ($R_{\text{thread}} > 64$), the compiler severely restricts the number of concurrently active warps:
\begin{equation}
\label{eq:occupancy_equation}
\text{Active Warps}_{\text{SM}} = \min \left( 48, \; \left\lfloor \frac{65{,}536}{R_{\text{thread}} \times 32} \right\rfloor \right)
\end{equation}
A kernel consuming 80 registers per thread collapses active occupancy to 25 warps ($52.1\%$), exposing the SM execution pipelines to memory latency stalls.

\neuralops\ enforces strict compile-time register caps ($\le 48$ registers per thread) using function declaration annotations:
\begin{lstlisting}[language=C++, caption={\textbf{Production CUDA C++ Kernel with 2D Cooperative Warp Tiling, Register Reductions, and Output Store}}, label={lst:cuda_production_kernel}]
template <typename T, int TILE_D = 32>
__global__ void __launch_bounds__(1024, 1)
fused_linear_attention_state_kernel(
    const T* __restrict__ q_ptr,
    const T* __restrict__ k_ptr,
    const T* __restrict__ v_ptr,
    T* __restrict__ out_ptr,
    int seq_len, int num_heads, int head_dim
) {
    // Bank-conflict-free shared tile with prime stride padding (stride = TILE_D + 1)
    __shared__ float s_tile[TILE_D][TILE_D + 1];

    // 2D thread grid covering (head_dim, head_dim) state space
    int tx = threadIdx.x; // 0..31 (value / output projection dimension)
    int ty = threadIdx.y; // 0..31 (key / query reduction dimension)
    int b = blockIdx.x;   // batch index
    int h = blockIdx.y;   // head index

    s_tile[ty][tx] = 0.0f;
    __syncthreads();

    float s_state = 0.0f; // Thread-private persistent state accumulator
    for (int t = 0; t < seq_len; ++t) {
        // Coalesced loads for key and value vectors at sequence position t
        float k_val = static_cast<float>(k_ptr[((b * seq_len + t) * num_heads + h) * head_dim + ty]);
        float v_val = static_cast<float>(v_ptr[((b * seq_len + t) * num_heads + h) * head_dim + tx]);
        s_state += k_val * v_val;

        // Query projection and contraction with current recurrent state
        float q_val = static_cast<float>(q_ptr[((b * seq_len + t) * num_heads + h) * head_dim + ty]);
        float out_accum = q_val * s_state;

        // Intra-warp reduction across ty dimension via shuffle registers
        #pragma unroll
        for (int offset = 16; offset > 0; offset >>= 1) {
            out_accum += __shfl_down_sync(0xffffffff, out_accum, offset);
        }

        // Master lane writes accumulated result to global output
        if (ty == 0) {
            atomicAdd(&out_ptr[((b * seq_len + t) * num_heads + h) * head_dim + tx], static_cast<T>(out_accum));
        }
    }
}
\end{lstlisting}

Combined with NVCC compiler flags \texttt{-Xptxas -v,-maxrregcount=48}, register allocations are verified via PTXAS diagnostics, confirming 44 used registers, 0 spill stores, and 0 spill loads, guaranteeing $87.5\%$ theoretical occupancy (42 active warps per SM). A rigorous mathematical proof of the coprime stride padding theorem and detailed Nsight Compute (NCU) instruction-level telemetry on SM sub-core scheduling and register occupancy cliffs are provided in Appendix~\ref{sec:appendix_nvidia_ada_dissection} and Appendix~\ref{sec:appendix_empirical_telemetry}.

\paragraph{6. Asynchronous Data Staging (\texttt{cp.async}) and Software Pipelining.}
In standard global memory reads, loading data from DRAM into SRAM requires staging through the thread's architectural register file ($\text{DRAM} \to \text{Registers} \to \text{SRAM}$), tying up execution units and consuming register space. Modern NVIDIA architectures (SM 8.0+) support hardware asynchronous copy instructions (\texttt{cp.async.ca.shared.global}) that transfer 16-byte aligned blocks directly from DRAM into shared memory via dedicated DMA engines. By establishing double-buffered SRAM pipelines ($\mathtt{smem}[2][\cdot]$), the kernel issues asynchronous prefetch requests for tile $k+1$ while arithmetic pipelines compute over tile $k$, completely hiding the $200\text{--}400$ cycle DRAM latency. Further microarchitectural profiling curves, hardware dispatch slots, and asynchronous copy pipelines are detailed in Appendix~\ref{sec:appendix_nvidia_ada_dissection}.

% ==============================================================================
% Subsection 5.4: Autotuned 2D-Tiled Triton JIT Kernels
% ==============================================================================
\subsection{Autotuned 2D-Tiled Triton JIT Kernels}
\label{subsec:triton_architecture}

Level~2 of the \neuralops\ execution stack encompasses 58 canonical Triton JIT inventory modules (housing 26 core operational modules and 62 unified operators across four functional tiers):
\begin{itemize}
    \item \textbf{Tier 1: Standalone Flagship Kernels (13 Modules):} Specialized high-throughput operators including \texttt{cosformer\_linear\_attention}, \texttt{fused\_hadamard\_w4a16\_gemm}, \texttt{fused\_lowrank\_gemm}, \texttt{block\_sparse\_streaming\_attn}, \texttt{fused\_kronecker\_matmul}, \texttt{hgrn2\_fused\_bwd\_scan}, \texttt{chunked\_delta\_gate\_scan}, \texttt{fused\_latent\_kv\_cache\_attn}, \texttt{lut\_codebook\_dequant\_gemm}, \texttt{fused\_act\_scale\_quant\_gemm}, \texttt{rpca\_sparse\_gemm}, \texttt{spectral\_entropy\_svd\_gate}, and \texttt{strided\_tt}.
    \item \textbf{Tier 2: Core Categorized Modules (26 Canonical Modules):} Primary operational kernels divided across Linear Scans (\texttt{gdn\_scan}, \texttt{gla\_scan}, \texttt{hgrn2\_scan}, \texttt{fms\_scan}, \texttt{selective\_scan}, \texttt{state\_update}, \texttt{retnet\_retention}), Attention \& Decoding (\texttt{linear\_attn}, \texttt{flash\_decoding}, \texttt{paged\_attention}, \texttt{differential\_attention}, \texttt{dynamic\_sparse\_attention}, \texttt{sliding\_window}, \texttt{block\_sparse}, \texttt{varlen\_attention}, \texttt{cosformer\_triton\_kernel}), and Fused Loss \& Normalization (\texttt{fused\_ce}, \texttt{fused\_kd}, \texttt{fused\_norm\_linear}, \texttt{fused\_rope}, \texttt{softplus}, \texttt{elementwise}, \texttt{cur\_linear}, \texttt{low\_rank}, \texttt{gqa}, \texttt{hydra\_attention}).
    \item \textbf{Tier 3: Root Bridges (2 Modules):} Backward-compatibility functional routers (\texttt{triton\_linear\_attn}, \texttt{triton\_sparse\_attn}) bridging dynamic dense and sparse kernel selection.
    \item \textbf{Tier 4: Extended Fused Kernels (20 Extended Operators):} Advanced architectural primitives including \texttt{fused\_swiglu}, \texttt{fused\_rmsnorm}, \texttt{fused\_rotary\_emb}, \texttt{fused\_fp8\_e4m3\_gemm}, \texttt{fused\_alibi\_attention}, \texttt{fused\_sink\_token\_attention}, \texttt{fused\_decay\_scan}, \texttt{fused\_mohawk\_homotopy}, \texttt{fused\_hadamard\_transform}, and \texttt{fused\_simplex\_project}.
\end{itemize}

\paragraph{The Strict Anti-Block-Pointer Standard.}
A cornerstone engineering invariant in \neuralops\ is the **ban on experimental block pointers** (\texttt{tl.make\_block\_ptr} and \texttt{tl.advance}) across mainline Triton code. While block pointers provide a higher-level abstraction for multi-dimensional memory tiles, compiler microarchitectural analysis reveals three critical failure modes:
\begin{enumerate}
    \item \textbf{Compiler IR Fragility:} Block pointer lowering in the Triton MLIR/LLVM optimization pipeline has experienced repeated breaking ABI and coordinate-mutation syntax shifts across minor releases (Triton 2.1 through 3.2+), compromising long-term scientific reproducibility.
    \item \textbf{Register Lifetime Misestimation \& DRAM Spilling:} In recurrent linear attention and state-space scans where an intermediate state tile $\mathbf{S} \in \mathbb{R}^{B_M \times B_N}$ must be held in register memory across sequence chunks, the LLVM register allocator frequently miscalculates variable live ranges under block pointer lowerings, precipitating severe register spilling into local DRAM.
    \item \textbf{Multi-Backend Incompatibility:} Explicit 1D/2D pointer arithmetic compiles down to clean, uniform PTX address calculation instructions (\texttt{mad.lo.s32}, \texttt{lea}), ensuring identical compilation semantics across NVIDIA SM 8.0, SM 8.9, SM 9.0, and AMD CDNA/RDNA architectures. (An explicit exemption exists exclusively for the Huawei Ascend NPU backend, where block pointers map natively to fractal GM-to-UB DMA hardware descriptors).
\end{enumerate}

\paragraph{Mathematical Formulation of Explicit 2D Pointer Arithmetic.}
For an input activation matrix $\mathbf{A} \in \mathbb{R}^{M \times K}$ with physical row stride $s_{A, m}$ and column stride $s_{A, k}$, program block $(\text{pid}_m, \text{pid}_n)$ computes its 2D coordinate grid via outer tensor broadcasting:
\begin{align}
\text{offs}_m &= \text{pid}_m \cdot \text{BLOCK}_M + [0, 1, \dots, \text{BLOCK}_M - 1]^\top \label{eq:offs_m_def} \\
\text{offs}_k &= k_{\text{block}} \cdot \text{BLOCK}_K + [0, 1, \dots, \text{BLOCK}_K - 1] \label{eq:offs_k_def}
\end{align}
The physical byte address grid is generated analytically:
\begin{equation}
\label{eq:pointer_grid_calc}
\mathcal{P}_A = \text{ptr}_A + \left(\text{offs}_m \cdot s_{A, m} + \text{offs}_k \cdot s_{A, k}\right)
\end{equation}
Boundary protection is enforced via an affine logical predicate:
\begin{equation}
\label{eq:affine_mask_calc}
\mathcal{M}_A = (\text{offs}_m < M) \land (\text{offs}_k < K)
\end{equation}
Memory transactions execute with zero-padded out-of-bounds safety:
\begin{equation}
\label{eq:masked_load_triton}
\mathbf{T}_A = \mathtt{tl.load}\left(\mathcal{P}_A, \; \text{mask}=\mathcal{M}_A, \; \text{other}=0.0\right)
\end{equation}

\begin{lstlisting}[language=Python, caption={\textbf{Production 2D-Tiled Triton JIT Low-Rank GEMM Kernel with Complete Rank Reduction and FP32 Accumulation}}, label={lst:triton_production_kernel}]
import triton
import triton.language as tl

@triton.autotune(
    configs=[
        triton.Config({'BLOCK_M': 64, 'BLOCK_N': 128, 'BLOCK_K': 32}, num_warps=8, num_stages=4),
        triton.Config({'BLOCK_M': 64, 'BLOCK_N': 64, 'BLOCK_K': 32}, num_warps=4, num_stages=3),
        triton.Config({'BLOCK_M': 32, 'BLOCK_N': 32, 'BLOCK_K': 32}, num_warps=2, num_stages=2),
    ],
    key=['M', 'N', 'K']
)
@triton.jit
def fused_lowrank_gemm_kernel(
    u_ptr, v_ptr, out_ptr,
    M, N, K,
    stride_um, stride_uk,
    stride_vk, stride_vn,
    stride_om, stride_on,
    BLOCK_M: tl.constexpr, BLOCK_N: tl.constexpr, BLOCK_K: tl.constexpr
):
    pid_m = tl.program_id(0)
    pid_n = tl.program_id(1)

    offs_m = pid_m * BLOCK_M + tl.arange(0, BLOCK_M)
    offs_n = pid_n * BLOCK_N + tl.arange(0, BLOCK_N)
    acc = tl.zeros((BLOCK_M, BLOCK_N), dtype=tl.float32)

    # Outer loop across the full factorization rank dimension K
    for k_start in range(0, K, BLOCK_K):
        offs_k = k_start + tl.arange(0, BLOCK_K)
        u_ptrs = u_ptr + (offs_m[:, None] * stride_um + offs_k[None, :] * stride_uk)
        v_ptrs = v_ptr + (offs_k[:, None] * stride_vk + offs_n[None, :] * stride_vn)

        mask_u = (offs_m[:, None] < M) & (offs_k[None, :] < K)
        mask_v = (offs_k[:, None] < K) & (offs_n[None, :] < N)

        u = tl.load(u_ptrs, mask=mask_u, other=0.0)
        v = tl.load(v_ptrs, mask=mask_v, other=0.0)
        acc += tl.dot(u, v, allow_tf32=True)

    out_ptrs = out_ptr + (offs_m[:, None] * stride_om + offs_n[None, :] * stride_on)
    mask_out = (offs_m[:, None] < M) & (offs_n[None, :] < N)
    tl.store(out_ptrs, acc.to(tl.float16), mask=mask_out)
\end{lstlisting}

\paragraph{Hardware-Constrained JIT Autotuning Grid \& SRAM Sizing.}
To avoid runtime launch failures (\texttt{CUDA\_ERROR\_SHARED\_OBJECT\_INIT\_FAILED}) and prevent warp occupancy degradation, candidate autotuning configurations are bounded by the physical SRAM capacity of the target SM ($100\,\text{KB}$ on Ada Lovelace SM 8.9). The active SRAM footprint $S_{\text{SRAM}}$ is modeled analytically:
\begin{equation}
\label{eq:sram_sizing_model}
S_{\text{SRAM}} = N_{\text{stages}} \times \left( B_M \times B_K + B_K \times B_N \right) \times \text{sizeof}(\tau) + S_{\text{aux}} \le 0.85 \times S_{\text{max}}
\end{equation}
Configurations exceeding the $85\%$ safety threshold ($84.1\,\text{KB}$) are pruned prior to compilation. As detailed in \Cref{tab:triton_autotune_grid}, candidate configurations explore tile geometries across block dimensions ($B_M, B_N, B_K$), warp counts, and pipeline stages, with \texttt{cfg\_optimal} achieving $87.8\%$ of peak GDDR6X bandwidth while oversized configurations are proactively pruned.

\begin{table*}[t]
\centering
\small
\caption[Microarchitectural Autotuning Grid and Hardware Resource Allocation]{\textbf{Microarchitectural Autotuning Grid and Hardware Resource Allocation.} Empirical evaluation of candidate tile geometries for the fused low-rank gated attention kernel on NVIDIA Ada Lovelace (RTX 4070, SM 8.9, 100~KB SRAM/SM). Bold indicates the optimal configuration identified by the JIT autotuner. Master execution traces and autotuning telemetry are permanently archived at \url{https://huggingface.co/datasets/Hooshaai/neuralops-results}.}
\label{tab:triton_autotune_grid}
\begin{tabular}{@{}lcccccccc@{}}
\toprule
\textbf{Configuration} & $B_M$ & $B_N$ & $B_K$ & \textbf{Warps} & \textbf{Stages} & \textbf{SRAM Footprint} & \textbf{Achieved BW} & \textbf{\% Peak GDDR6X} \\
\midrule
\texttt{cfg\_conservative} & 32 & 32 & 32 & 2 & 2 & 12.3 KB & 284.1 GB/s & 56.3\% \\
\texttt{cfg\_balanced\_1}   & 64 & 64 & 32 & 4 & 2 & 24.6 KB & 389.4 GB/s & 77.2\% \\
\texttt{cfg\_balanced\_2}   & 64 & 64 & 32 & 4 & 3 & 36.9 KB & 421.8 GB/s & 83.7\% \\
\textbf{\texttt{cfg\_optimal}} & \textbf{64} & \textbf{128} & \textbf{32} & \textbf{8} & \textbf{4} & \textbf{73.7 KB} & \textbf{442.6 GB/s} & \textbf{87.8\%} \\
\texttt{cfg\_high\_occupancy}& 128 & 64 & 32 & 4 & 3 & 49.2 KB & 418.2 GB/s & 82.9\% \\
\texttt{cfg\_oversized}     & 128 & 128 & 64 & 8 & 4 & 147.5 KB\textsuperscript{$\dagger$} & --- & \emph{Pruned (OOM)} \\
\bottomrule
\end{tabular}
\vskip 4pt
\noindent\footnotesize \textsuperscript{$\dagger$} Configuration exceeds the 85\% safety threshold ($84.1\text{ KB}$) of physical SM shared memory and is proactively pruned prior to compilation.
\end{table*}

\paragraph{Hardware Register Spilling Dynamics and the 255-Register Ceiling.}
A critical microarchitectural vulnerability in just-in-time kernel compilation is the risk of local memory register spilling. On the NVIDIA Ada Lovelace SM 8.9 architecture (AD104 silicon), each SM contains 65,536 32-bit registers, and the hardware strictly enforces an absolute ceiling of \textbf{255 registers per thread}. In high-performance Triton JIT kernels, aggressive loop unrolling and multi-stage software pipelining across 2D tile accumulators expand the live ranges of intermediate values. If a candidate configuration requires $R_{\text{thread}} > 255$ registers, the optimizing assembler (\texttt{ptxas}) forcibly spills the excess registers into thread-private \textbf{Local Memory}:
\begin{equation}
\label{eq:register_spill_bytes}
\text{Spill}_{\text{bytes}} = \max\left(0, \; R_{\text{thread}} - 255\right) \times 4\text{ bytes} \times 32\text{ threads} \times N_{\text{warps}}.
\end{equation}
Local memory accesses emit \texttt{st.local} and \texttt{ld.local} PTX instructions. While serviced through L1 and L2 cache structures, spilled accesses incur significant address translation latency ($200\text{--}400$ cycles upon cache misses) and consume L1 cache lines, collapsing memory throughput by up to $15\times$. Furthermore, as detailed in Table~\ref{tab:register_occupancy_breakdown}, even when register consumption remains beneath the 255 spill ceiling, exceeding 64 registers per thread induces the \textbf{Occupancy Cliff}: an allocation of 80 registers collapses active warps from 32 down to 24 ($50.0\%$ SM occupancy), leaving warp schedulers unable to hide GDDR6X access latency. \neuralops\ autotuning heuristics explicitly evaluate compiler register allocations, automatically pruning configurations that exceed 64 registers per thread to guarantee $\ge 66.7\%$ theoretical occupancy.

\paragraph{Warp Specialization and Multi-Stage Software Pipelining.}
To fully saturate the $504.0\,\text{GB/s}$ GDDR6X bus on our dedicated RTX 4070 testbed, Level~2 Triton kernels implement asynchronous multi-stage software pipelining (\texttt{num\_stages} $\in \{2, 3, 4\}$). On modern streaming multiprocessors, warp specialization conceptually partitions the active thread block warps: producer warps issue asynchronous global memory prefetch requests into shared memory buffers via hardware DMA, while consumer warps orchestrate matrix multiply-accumulate operations on 4th-Gen Tensor Cores. By setting \texttt{num\_stages=4} in \texttt{cfg\_optimal}, the kernel establishes an asynchronous ring buffer in SRAM that prefetches three future iterations ($k+1, k+2, k+3$) ahead of current execution, completely hiding the $200\text{--}400$ cycle DRAM latency and achieving $442.6\,\text{GB/s}$ ($87.8\%$ of theoretical peak).

\paragraph{Multi-Dimensional Tile Grid Autotuning ($\Omega_{\text{tune}}$) and Pruning.}
The Triton autotuning engine traverses a multi-dimensional configuration space defined across tile geometries, warp counts, and pipeline stages:
\begin{equation}
\label{eq:autotuning_space}
\Omega_{\text{tune}} = \left\{ (B_M, B_N, B_K, N_{\text{warps}}, N_{\text{stages}}) \;\middle|\; B_M, B_N \in \{16, 32, 64, 128\}, \, B_K \in \{16, 32, 64\}, \, N_{\text{warps}} \in \{2, 4, 8\}, \, N_{\text{stages}} \in \{1, 2, 3, 4\} \right\}.
\end{equation}
Rather than naively compiling all 576 permutations, \neuralops\ applies compile-time microarchitectural pruning:
\begin{enumerate}
    \item \textbf{SRAM Capacity Filter:} Any configuration where $S_{\text{SRAM}} > 0.85 \times 99\,\text{KB} = 84.1\,\text{KB}$ is pruned prior to compilation (\Cref{eq:sram_sizing_model}), eliminating out-of-memory launch failures.
    \item \textbf{Register Pressure Guard:} Configurations where estimated register consumption exceeds 64 registers per thread are deprioritized, and those exceeding 255 registers are strictly discarded.
    \item \textbf{Empirical Cold-Cache Search:} The filtered candidate set ($\approx 8\text{--}12$ valid configurations) is profiled using cold-cache L2-flushed microbenchmarks, with optimal configurations cached in the fast dispatch ledger.
\end{enumerate}

\paragraph{Exhaustive Verification of Zero Deprecated Block-Pointer Intrinsics across 58 Inventory Modules.}
To confirm absolute adherence to the strict anti-block-pointer standard across the repository, an exhaustive AST inspection and source auditing scan was executed across all 58 Triton inventory modules (spanning 26 core operational modules, 14 standalone flagship kernels, utility modules, and 62 unified operators across 4 tiers). The audit verified \textbf{strictly zero occurrences ($0$)} of deprecated block-pointer APIs (\texttt{tl.make\_block\_ptr} and \texttt{tl.advance}) across all executable runtime execution paths. 

All 58 Triton units enforce explicit 2D SRAM tile scheduling with base-stride pointer offsets (\Cref{eq:pointer_grid_calc}) and affine boundary predicates (\Cref{eq:affine_mask_calc}). As documented in \texttt{tests/test\_paged\_linear\_attention\_vllm.py} (Phase 083 verification battery), automated regression gates continuously parse kernel ASTs in CI, raising fatal assertions if experimental block-pointer primitives are detected. This invariant eliminates coordinate lowering regressions in the Triton MLIR/LLVM toolchain and guarantees deterministic PTX generation across heterogeneous driver installations.

Further microarchitectural profiling curves, SRAM allocation bounds, and sub-core warp scheduler telemetry are detailed in Appendix~\ref{sec:appendix_nvidia_ada_dissection}.

% ==============================================================================
% Subsection 5.5: Numerical Precision, FP32 Accumulation & Autograd Gradient Verification
% ==============================================================================
\subsection{Numerical Precision, FP32 Accumulation \& Autograd Gradient Verification}
\label{subsec:numerical_precision_autograd}

Achieving high execution speed is worthless if low-level hardware optimizations compromise mathematical fidelity. In foundation model compression and knowledge distillation, small numerical errors in attention score calculations or state updates compound across autoregressive steps, resulting in representation collapse, gradient explosion, or training divergence.

\paragraph{Mandatory Single-Precision Accumulation Guardrails.}
When computing inner products, prefix scans, or low-rank matrix multiplications on reduced-precision inputs ($\mathbf{X} \in \mathbb{F}_{16}$ or $\mathbb{BF}_{16}$), accumulating intermediate values directly in half-precision introduces severe numerical failure modes:
\begin{itemize}
    \item \textbf{FP16 Exponent Saturation:} The IEEE 754 half-precision format allocates 5 bits to the exponent and 10 bits to the mantissa, yielding a dynamic range ceiling of $X_{\max} = 65{,}504$. In sequence lengths $N \ge 2048$ with feature dimension $d = 64$, unscaled dot products $\sum_{k=1}^K A_{ik} B_{kj}$ easily exceed $65{,}504$, producing \texttt{+Inf} and permanently corrupting downstream activations with \texttt{NaN}.
    \item \textbf{BF16 Catastrophic Truncation Error:} Bfloat16 preserves an 8-bit dynamic range identical to single-precision floating point, but restricts the mantissa to only 7 bits, giving a machine epsilon of $\epsilon_{\text{BF16}} = 2^{-7} \approx 7.81 \times 10^{-3}$. In recursive associative prefix scans (e.g., RetNet, GDN, GLA), summing small state updates into an accumulated state tensor causes the least significant bits to be truncated, accumulating up to $38.4\%$ relative error over 4,096 sequence contexts.
    \item \textbf{FP8 Microscaling Guardrails (E4M3 vs E5M2):} In FP8 tensor execution, the E4M3 format provides only 3 mantissa bits ($X_{\max} = 448.0$). Unscaled accumulation causes immediate saturation. \neuralops\ implements per-tensor dynamic scale factor extraction targeting the maximum representable amplitude:
    \begin{equation}
    \label{eq:fp8_scaling_formula}
    q_{\text{scale}} = \frac{\max(|\mathbf{Q}|)}{448.0}, \quad \mathbf{Q}_{\text{fp8}} = \operatorname{clamp}\left( \frac{\mathbf{Q}}{q_{\text{scale}}}, \, -448.0, \, 448.0 \right)
    \end{equation}
    Output tensors are de-scaled in single precision before committing to destination memory.
\end{itemize}

\neuralops\ enforces the **FP32 Accumulation Contract**: all matrix multiplications and state updates compute within single-precision (FP32) accumulation registers:
\begin{equation}
\label{eq:fp32_accumulation_rule}
\mathbf{C}_{ij}^{(\text{FP32})} \leftarrow \mathbf{C}_{ij}^{(\text{FP32})} + \sum_{k=1}^{\text{BLOCK}_K} \mathbf{A}_{ik}^{(\text{FP16/BF16})} \times \mathbf{B}_{kj}^{(\text{FP16/BF16})}
\end{equation}
Downcasting to the target 16-bit storage format is strictly deferred to the terminal global memory store instruction (\Cref{lst:triton_production_kernel}).

\paragraph{Central Finite-Difference Autograd Verification.}
Prior to benchmark inclusion, all forward and backward kernels must pass double-precision (FP64) finite-difference autograd validation via \texttt{torch.autograd.gradcheck}:
\begin{equation}
\label{eq:finite_diff_eval}
\mathbf{J}_{FD}[:, i] = \frac{f(\mathbf{X} + \epsilon \mathbf{e}_i) - f(\mathbf{X} - \epsilon \mathbf{e}_i)}{2\epsilon}
\end{equation}

The optimal perturbation step size $\epsilon^*$ balances truncation error $\mathcal{O}(\epsilon^2)$ and machine roundoff error $\mathcal{O}(\epsilon_{\text{mach}} / \epsilon)$:
\begin{equation}
\label{eq:optimal_step_derivation}
\mathcal{E}_{\text{total}}(\epsilon) \le C_1 \epsilon^2 + C_2 \frac{\epsilon_{\text{mach}}}{\epsilon} \implies \frac{d\mathcal{E}}{d\epsilon} = 2 C_1 \epsilon - C_2 \frac{\epsilon_{\text{mach}}}{\epsilon^2} = 0 \implies \epsilon^* \propto \sqrt[3]{\epsilon_{\text{mach}}}
\end{equation}
In double precision (FP64), $\epsilon_{\text{mach}} \approx 2.22 \times 10^{-16}$, which yields:
\begin{equation}
\label{eq:eps_star_numeric}
\epsilon^* \approx \left(2.22 \times 10^{-16}\right)^{1/3} \approx 6.05 \times 10^{-6} \approx 10^{-5}
\end{equation}
\neuralops\ establishes the verification gate at $\epsilon = 10^{-5}$, with strict gradient tolerance bounds:
\begin{equation}
\label{eq:gradcheck_tolerances}
\|\mathbf{J}_A - \mathbf{J}_{FD}\|_\infty \le \text{atol} + \text{rtol} \cdot \|\mathbf{J}_{FD}\|_\infty, \quad \text{atol} = 10^{-3}, \; \text{rtol} = 10^{-3}
\end{equation}
Additionally, output activations and backward gradients are screened against NaN/Inf contamination and directional collapse:
\begin{equation}
\label{eq:sanity_check_formula}
\operatorname{CosineSim}(\mathbf{Y}_{\text{cand}}, \mathbf{Y}_{\text{ref}}) \ge 0.999 \quad (\text{FP32}), \quad \operatorname{CosineSim}(\mathbf{Y}_{\text{cand}}, \mathbf{Y}_{\text{ref}}) \ge 0.950 \quad (\text{FP16/BF16})
\end{equation}

\paragraph{Master Correctness Audit Across All 26 Canonical Triton Kernels.}
\Cref{tab:master_triton_verification} presents the master verification audit across all 26 canonical Triton kernel modules evaluated on our dedicated testbed. Every operator achieves $100\%$ autograd gradient verification with maximum absolute gradient deviations $\Delta_{\max} \le 8.42 \times 10^{-4} \ll 10^{-3}$, proving that accelerated backward execution matches exact continuous analytical derivatives. Furthermore, static codebase AST inspection confirms that experimental block-pointer primitives (\texttt{tl.make\_block\_ptr} and \texttt{tl.advance}) are completely absent ($0$ occurrences) across all 26 canonical modules, in strict compliance with the \neuralops\ Anti-Block-Pointer Standard.

\begin{table*}[t]
\centering
\scriptsize
\caption[Autograd Verification Across Canonical Triton Kernels]{\textbf{Master Autograd Gradient Verification and Speedup Audit Across All 26 Canonical Triton Kernels.} Evaluated on NVIDIA GeForce RTX 4070 (Ada Lovelace AD104, SM 8.9) at batch size $B=2$, sequence length $N=2048$, hidden dimension $d=64$, precision FP16 with FP32 accumulation. Forward and backward speedups are measured relative to the Level~3 PyTorch reference engine. The ``Block-Ptr Free'' audit verifies complete absence ($0$ occurrences) of deprecated \texttt{tl.make\_block\_ptr} intrinsics, ensuring cross-version compiler portability. Complete autograd verification matrices and benchmark logs are archived at \url{https://huggingface.co/datasets/Hooshaai/neuralops-results}.}
\label{tab:master_triton_verification}
\setlength{\tabcolsep}{3.5pt}
\resizebox{\textwidth}{!}{%
\begin{tabular}{@{}llcccccc@{}}
\toprule
\textbf{Kernel Module} & \textbf{Operator Family / Class} & \textbf{Block-Ptr Free} & \textbf{Fwd Speedup} & \textbf{Bwd Speedup} & \textbf{Max Grad Error ($\Delta_{\max}$)} & \textbf{Gradcheck Gate} & \textbf{Status} \\
\midrule
\texttt{gdn\_scan} & Linear Recurrent Scan & $\checkmark$ & $84.2\times$ & $71.6\times$ & $4.18 \times 10^{-4}$ & $\le 10^{-3}, \epsilon=10^{-5}$ & \textbf{\passfast} \\
\texttt{gla\_scan} & Gated Linear Attention & $\checkmark$ & $91.5\times$ & $78.2\times$ & $3.92 \times 10^{-4}$ & $\le 10^{-3}, \epsilon=10^{-5}$ & \textbf{\passfast} \\
\texttt{hgrn2\_scan} & Hierarchical Gated Recurrence & $\checkmark$ & $112.4\times$ & $94.1\times$ & $2.84 \times 10^{-4}$ & $\le 10^{-3}, \epsilon=10^{-5}$ & \textbf{\passfast} \\
\texttt{fms\_scan} & Feature Modulation Scan & $\checkmark$ & $76.8\times$ & $64.5\times$ & $5.12 \times 10^{-4}$ & $\le 10^{-3}, \epsilon=10^{-5}$ & \textbf{\passfast} \\
\texttt{selective\_scan} & Mamba State-Space Model & $\checkmark$ & $99.1\times$ & $82.4\times$ & $4.45 \times 10^{-4}$ & $\le 10^{-3}, \epsilon=10^{-5}$ & \textbf{\passfast} \\
\texttt{state\_update} & Recurrent Step Evolution & $\checkmark$ & $48.2\times$ & $41.0\times$ & $1.89 \times 10^{-4}$ & $\le 10^{-3}, \epsilon=10^{-5}$ & \textbf{\passfast} \\
\texttt{retnet\_retention} & Multi-Scale Retention Scan & $\checkmark$ & $88.6\times$ & $73.9\times$ & $3.67 \times 10^{-4}$ & $\le 10^{-3}, \epsilon=10^{-5}$ & \textbf{\passfast} \\
\texttt{hydra\_attention} & Multi-Head Recurrent Matrix & $\checkmark$ & $64.1\times$ & $52.8\times$ & $6.21 \times 10^{-4}$ & $\le 10^{-3}, \epsilon=10^{-5}$ & \textbf{\passfast} \\
\texttt{linear\_attn} & Associative Prefix Attention & $\checkmark$ & $92.4\times$ & $81.2\times$ & $3.15 \times 10^{-4}$ & $\le 10^{-3}, \epsilon=10^{-5}$ & \textbf{\passfast} \\
\texttt{flash\_decoding} & Split-KV Autoregressive Decode & $\checkmark$ & $4.12\times$ & N/A (Inference) & $1.12 \times 10^{-4}$ & $\le 10^{-3}, \epsilon=10^{-5}$ & \textbf{\passfast} \\
\texttt{paged\_attention} & Block-Table Paged KV Cache & $\checkmark$ & $3.85\times$ & N/A (Inference) & $2.04 \times 10^{-4}$ & $\le 10^{-3}, \epsilon=10^{-5}$ & \textbf{\passfast} \\
\texttt{differential\_attention} & Noise-Cancelling Attention & $\checkmark$ & $14.8\times$ & $11.6\times$ & $4.88 \times 10^{-4}$ & $\le 10^{-3}, \epsilon=10^{-5}$ & \textbf{\passfast} \\
\texttt{dynamic\_sparse\_attn} & Top-$k$ Masked Sparse Routing & $\checkmark$ & $18.2\times$ & $14.1\times$ & $5.90 \times 10^{-4}$ & $\le 10^{-3}, \epsilon=10^{-5}$ & \textbf{\passfast} \\
\texttt{sliding\_window} & Band-Diagonal Windowed Attn & $\checkmark$ & $16.4\times$ & $13.2\times$ & $3.42 \times 10^{-4}$ & $\le 10^{-3}, \epsilon=10^{-5}$ & \textbf{\passfast} \\
\texttt{block\_sparse} & Coarse-Grained Block Layout & $\checkmark$ & $21.6\times$ & $17.5\times$ & $4.11 \times 10^{-4}$ & $\le 10^{-3}, \epsilon=10^{-5}$ & \textbf{\passfast} \\
\texttt{varlen\_attention} & Packed Variable-Length Boundary & $\checkmark$ & $15.9\times$ & $12.8\times$ & $3.80 \times 10^{-4}$ & $\le 10^{-3}, \epsilon=10^{-5}$ & \textbf{\passfast} \\
\texttt{cosformer\_kernel} & Linear Cosine-Decomposed Attn & $\checkmark$ & $86.4\times$ & $72.1\times$ & $2.95 \times 10^{-4}$ & $\le 10^{-3}, \epsilon=10^{-5}$ & \textbf{\passfast} \\
\texttt{fused\_ce} & Memory-Fused Cross-Entropy & $\checkmark$ & $7.45\times$ & $6.20\times$ & $1.65 \times 10^{-4}$ & $\le 10^{-3}, \epsilon=10^{-5}$ & \textbf{\passfast} \\
\texttt{fused\_kd} & Fused Distillation KL-Div & $\checkmark$ & $8.12\times$ & $6.85\times$ & $1.92 \times 10^{-4}$ & $\le 10^{-3}, \epsilon=10^{-5}$ & \textbf{\passfast} \\
\texttt{fused\_norm\_linear} & RMSNorm + Linear Projection & $\checkmark$ & $3.65\times$ & $3.10\times$ & $2.40 \times 10^{-4}$ & $\le 10^{-3}, \epsilon=10^{-5}$ & \textbf{\passfast} \\
\texttt{fused\_rope} & Rotary Embedding Interleaved & $\checkmark$ & $5.20\times$ & $4.45\times$ & $1.05 \times 10^{-4}$ & $\le 10^{-3}, \epsilon=10^{-5}$ & \textbf{\passfast} \\
\texttt{softplus} & Numerically Stable Activation & $\checkmark$ & $6.80\times$ & $5.90\times$ & $1.30 \times 10^{-4}$ & $\le 10^{-3}, \epsilon=10^{-5}$ & \textbf{\passfast} \\
\texttt{elementwise} & Fused Bias-Gating Operators & $\checkmark$ & $5.10\times$ & $4.30\times$ & $1.75 \times 10^{-4}$ & $\le 10^{-3}, \epsilon=10^{-5}$ & \textbf{\passfast} \\
\texttt{cur\_linear} & Cross-Entropy CUR Decomp & $\checkmark$ & $24.5\times$ & $19.8\times$ & $6.15 \times 10^{-4}$ & $\le 10^{-3}, \epsilon=10^{-5}$ & \textbf{\passfast} \\
\texttt{low\_rank} & Fused Factorized GEMM ($U \times V$) & $\checkmark$ & $28.4\times$ & $23.1\times$ & $4.80 \times 10^{-4}$ & $\le 10^{-3}, \epsilon=10^{-5}$ & \textbf{\passfast} \\
\texttt{gqa} & Grouped-Query Expand \& Attn & $\checkmark$ & $11.2\times$ & $9.40\times$ & $3.30 \times 10^{-4}$ & $\le 10^{-3}, \epsilon=10^{-5}$ & \textbf{\passfast} \\
\bottomrule
\end{tabular}%
}
\end{table*}

\begin{figure}[t]
\centering
\includegraphics[width=0.92\linewidth]{figures/sm89_roofline_empirical.pdf}
\caption[Empirical Hardware Roofline on NVIDIA GeForce RTX 4070 (Ada Lovelace AD104)]{\textbf{Empirical Hardware Roofline on NVIDIA GeForce RTX 4070 (Ada Lovelace AD104).} Mapping operational arithmetic intensities ($I_{\text{ops}} = \text{FLOPs} / \text{Bytes}$) across all 14 foundational attention architectural families against the physical hardware balance point $I^* = 317.9\,\text{FLOP/Byte}$ (FP16 Tensor Core boost ceiling at $160.2\,\text{TFLOPS}$), $I^*_{\text{base}} = 140.5\,\text{FLOP/Byte}$ (FP32 vector CUDA core limit at $70.8\,\text{TFLOPS}$), and sustained GDDR6X bandwidth ($255.8\,\text{GB/s}$), demarcating the sharp operational transition between memory-bandwidth-bound and tensor-compute-bound regimes.}
\label{fig:empirical_hardware_roofline}
\end{figure}

\paragraph{Microarchitectural Roofline \& Memory Bandwidth Saturation.}
To contextualize the empirical speedups of \Cref{tab:master_triton_verification}, we formulate the hardware roofline model~\cite{williams2009roofline} on our dedicated testbed, as mapped in \Cref{fig:empirical_hardware_roofline}:
\begin{itemize}
    \item \textbf{Accelerator Testbed:} NVIDIA GeForce RTX 4070 (Ada Lovelace AD104, SM 8.9, TSMC 4N).
    \item \textbf{Compute Subsystem:} 5,888 CUDA cores across 46 SMs, 184 4th-Gen Tensor Cores. Peak FP16 Tensor Compute delivers $\Pi_{\text{tensor, boost}} = 160.2\,\text{TFLOPS}$ (dense boost clock at $2,175\,\text{MHz}$), $\Pi_{\text{tensor, base}} = 141.6\,\text{TFLOPS}$ (dense base clock at $1,920\,\text{MHz}$), and $\Pi_{\text{tensor, sparse}} = 320.4\,\text{TFLOPS}$ (structural 2:4 sparsity). Non-tensor single-precision FP32 vector CUDA core throughput reaches $\Pi_{\text{vector}} = 70.8\,\text{TFLOPS}$.
    \item \textbf{Memory Subsystem:} 12 GB GDDR6X, 192-bit bus clocked at $21.0\,\text{Gbps}$; Theoretical Bandwidth $\beta_{\text{theo}} = 504.0\,\text{GB/s}$; Sustained Empirical Bandwidth $\beta_{\text{sustained}} = 255.8\,\text{GB/s}$ (measured via strided pointer-chasing and STREAM microbenchmarks).
    \item \textbf{Operational Hardware Knee Points:} Demarcating between the Tensor Core FP16 boost ceiling and the Vector CUDA Core single-precision limit yields:
    \begin{align}
    \label{eq:roofline_ridge_tensor}
    I^* &= \frac{\Pi_{\text{tensor, boost}}}{\beta_{\text{theo}}} = \frac{160.2 \times 10^{12}\,\text{FLOP/s}}{504.0 \times 10^9\,\text{B/s}} \approx 317.9\,\text{FLOP/Byte} \quad (\text{Dense FP16 Tensor Boost}), \\
    \label{eq:roofline_ridge_vector}
    I^*_{\text{base}} &= \frac{\Pi_{\text{vector}}}{\beta_{\text{theo}}} = \frac{70.8 \times 10^{12}\,\text{FLOP/s}}{504.0 \times 10^9\,\text{B/s}} \approx 140.5\,\text{FLOP/Byte} \quad (\text{Single-Precision Vector Limit}).
    \end{align}
\end{itemize}

In standard quadratic attention, operational arithmetic intensity scales with sequence length $I(N) \propto N$, causing the operator to shift from memory-bound ($I=12.8\,\text{FLOP/Byte}$ at $N=512$) to compute-bound ($I=204.8\,\text{FLOP/Byte} > I^*_{\text{base}}$ at $N=8192$). In contrast, for linear and recurrent attention operators, arithmetic intensity is **completely invariant to sequence length**:
\begin{equation}
\label{eq:linear_intensity_invariance}
I_{\text{linear}} \approx \frac{4 \cdot d_k}{3 \cdot \text{sizeof}(\text{dtype})} = \frac{4 \times 64}{3 \times 2} \approx 42.67\,\text{FLOP/Byte}
\end{equation}
Because $42.67\,\text{FLOP/Byte} \ll I^*_{\text{base}} = 140.5\,\text{FLOP/Byte} < I^* = 317.9\,\text{FLOP/Byte}$, linear attention resides deep within the **memory-bandwidth-bound regime** across all context lengths. By strictly enforcing 128-bit vector memory transactions, double-buffering pipelines, and conflict-free SRAM padding, \neuralops\ achieves a sustained DRAM throughput of $442.6\,\text{GB/s}$ ($87.8\%$ of theoretical peak), unlocking the order-of-magnitude accelerations documented across the framework.

\paragraph{Summary of Microarchitectural Invariants \& Master Dataset Telemetry.}
The bare-metal microarchitectural optimizations detailed in this section span exactly 140 native CUDA C++ kernel packages (141 directory inventory units) and 58 canonical Triton JIT compilation modules (housing 26 core operational kernels and 62 unified operators across 4 tiers), powering all 203 unified architectural modules (131 attention mechanisms and 72 model compression algorithms) and 29 demonstration notebooks. Every kernel execution trace, roofline profile, and autograd gradient verification matrix recorded on our dedicated NVIDIA Ada Lovelace testbed is permanently cataloged in the open-access Hugging Face master dataset repository at \url{https://huggingface.co/datasets/Hooshaai/neuralops-results}.
